\section{Graphics}
\label{sec:graphics}

Three kinds of figure are supported, and they have different rendering
stories, so pick deliberately:

\begin{itemize}[nosep]
  \item \textbf{PSTricks} (\cref{fig:pie,fig:plot}) --- rendered by
    \texttt{dvips}/Ghostscript in the PDF build. The subset used here
    (\texttt{pspicture}, \texttt{psline}, \texttt{pscircle}, \texttt{pswedge},
    \texttt{psplot}, \texttt{psaxes}, \texttt{rput}, \texttt{psdots}) is also
    what LaTeX2JS renders in the browser, so these figures can be reused as-is
    on Mathapedia.
  \item \textbf{TikZ} (\cref{fig:flow}) --- rendered in the PDF build only.
    Use it for diagrams that will never need a browser rendering.
  \item \textbf{Generated EPS} (\cref{fig:pipeline}) --- anything produced by an
    external tool (Mermaid here). \texttt{make figures} turns
    \texttt{tex/figures/*.mmd} into \texttt{build/figures/*.eps}; only the
    source is committed.
\end{itemize}

\begin{figure}[ht]
  \centering
  % PSTricks: circle, wedges, labels, arrows. Stays inside the LaTeX2JS subset
% so the same source renders on Mathapedia. Keep PSTricks commands at column
% 0 inside pspicture: the browser parser does not tolerate leading whitespace.
\begin{pspicture}(-3.5,-3.5)(5,3.5)
\pscircle[linewidth=0.5pt](0,0){3}
\pswedge[fillstyle=solid,fillcolor=figaccent,linestyle=none](0,0){3}{60}{120}
\pswedge[fillstyle=solid,fillcolor=figprimary,linestyle=none](0,0){3}{120}{420}
\rput(4,2.5){\textbf{one sixth}}
\psline{->}(3.2,2.3)(0.3,2.3)
\rput(-2.6,-2.8){\textbf{five sixths}}
\psline{->}(-2.4,-2.4)(-1.3,-1.3)
\psdots[dotstyle=o,dotsize=3pt](0,0)
\end{pspicture}

  \caption{A PSTricks figure. Everything in it is in the LaTeX2JS subset.}
  \label{fig:pie}
\end{figure}

\begin{figure}[ht]
  \centering
  % PSTricks plot (pst-plot). psaxes + psplot are in the LaTeX2JS subset.
\begin{pspicture}(-0.5,-1.5)(6.5,1.7)
\psaxes[labels=none,ticks=none]{->}(0,0)(-0.5,-1.5)(6.5,1.5)
\psplot[linecolor=figprimary,linewidth=1.2pt,plotpoints=200]{0}{6.28}{x 57.2958 mul sin}
\psplot[linecolor=figaccent,linewidth=1.2pt,plotpoints=200]{0}{6.28}{x 57.2958 mul cos}
\rput(6.5,-0.3){$x$}
\rput(-0.3,1.5){$y$}
\rput(3.2,1.4){\color{figprimary}$\sin x$}
\rput(1.4,-1.3){\color{figaccent}$\cos x$}
\end{pspicture}

  \caption{A PSTricks plot with axes; \texttt{psplot} is interactive on Mathapedia.}
  \label{fig:plot}
\end{figure}

\begin{figure}[ht]
  \centering
  % TikZ: the build pipeline as a flow diagram. PDF only (no browser rendering).
\begin{tikzpicture}[
  node distance=1.1cm and 0.9cm,
  box/.style={draw,rounded corners,minimum width=2cm,minimum height=0.8cm,
              align=center,font=\small},
  src/.style={box,fill=figprimary!15},
  art/.style={box,fill=figaccent!20},
  >={Stealth[length=2.2mm]}
]
\node[src] (tex) {\texttt{tex/*.tex}};
\node[src,below=of tex] (mmd) {\texttt{figures/*.mmd}};
\node[art,right=of mmd] (eps) {\texttt{build/figures/*.eps}};
\node[art,right=of tex] (dvi) {\texttt{main.dvi}};
\node[art,right=of dvi] (ps) {\texttt{main.ps}};
\node[art,right=of ps] (pdf) {\texttt{main.pdf}};
\node[art,below=of pdf] (svg) {\texttt{svg/*.svg}};

\draw[->] (tex) -- node[above,font=\scriptsize] {latex} (dvi);
\draw[->] (dvi) -- node[above,font=\scriptsize] {dvips} (ps);
\draw[->] (ps) -- node[above,font=\scriptsize] {ps2pdf} (pdf);
\draw[->] (mmd) -- node[above,font=\scriptsize] {mmdc + gs} (eps);
\draw[->] (eps) -| (dvi);
\draw[->] (pdf) -- node[right,font=\scriptsize] {dvisvgm} (svg);
\end{tikzpicture}

  \caption{A TikZ figure. PDF only.}
  \label{fig:flow}
\end{figure}

\begin{figure}[ht]
  \centering
  \includegraphics[width=0.85\linewidth]{pipeline}
  \caption{A Mermaid diagram compiled to EPS by \texttt{make figures}.}
  \label{fig:pipeline}
\end{figure}
